\section{研究課題}
GTFSを利用して，バスサイネージシステムに必要な情報を自動的に取り出す際に生じる問題を整理する．
GTFSでは，運行情報そのものは構造化されている一方で，バス利用者向けの意味情報がすべて構造化されているわけではない．

\subsection{headsignの構造化}
GTFSには，便の行先を表すtrip\_headsign・stop\_headsignというフィールドが存在する（以降2つをまとめて``headsign''とする）．
trip\_headsignは便全体の基本となる行先表示であり，stop\_headsignは特定の停留所でそれを上書き表示するものである．
GTFS-JPでは，車両の方向幕の表示を設定することを原則とし，急行、直通等の種別がある場合は行先に加えて種別を併記することになっている~\cite{GTFSJPv4}．
このフィールドには方向幕に書かれている文字列がそのまま書き込まれ，主に行先・経由・種別の3情報が含まれる．
これらの情報は，バスサイネージシステムの設定を最適化する上で有益であるため，構造化して取り出したい．
しかしGTFS仕様ではheadsign文字列の内部構造は規定されておらず，事業者によって表記方法が異なる．

% TODO: 表記例を入れる

これらをルールベースで解析する場合，いくつかの問題がある．
まず，事業者ごとに解析ルールを用意する必要がある．
事業者ごとにフォーマットが異なるため，内容を精査した上で，適切な抽出条件を事業者別に作成する必要がある．
GTFSを公開する全事業者に対して抽出条件を事前に用意することは現実的でない．
設置者に条件設定をさせることも，「設置者自身による設定」を目指す本システムでは採用できない．
また，事業者内でも複数フォーマットが存在している場合がある．
そのような場合，解析ルールは更に複雑なものになる．
加えて，停留所名に見える語彙があっても，それが直接駅名を表しているとは限らない．例えば函館バスの場合，「函館高専」「高専」と表記されている場合は，「函館高専前」という停留所を示している．青森市営バスの場合，「県病」と表記されている場合，「県立中央病院」という停留所を示している．他にも，函館バスでは「鍛治線」「桐花通」のように通りの名前が記述されている場合もある．
ここまでの整理から，必要なのは文字列を単純に分割する処理ではなく、headsignが表す意味を解釈する処理であるといえる．
そこで，AIエージェントによって内容を解釈させ，構造化されていないheadsignの情報を構造化する．

\subsection{方面別分類}

バスサイネージシステムにおいては、便を利用者が理解しやすい「方面」に分類することが必要である．
しかし，GTFSには経路の情報が存在するだけで，あるバス停から見た「方面」という情報は直接書き込まれているわけではない．
判断や表現に利用できる情報としては以下がある．

% textlint-disable
\begin{itemize}
    \item trip\_headsign，stop\_headsignに書き込まれている内容（行先・代表的な経由地）
    \item stop\_times.txtに書き込まれている停車駅の情報
    \item stops.txtに書き込まれている駅名などの，事業者が案内に用いている語彙
\end{itemize}
% textlint-enable

必要な情報がそろったとしても，単純なルールで方面を分類するのは難しい．
事業者や利用者が認識している「方面」には明確な判断基準は存在せず，感覚によって決められている．
路線網や地域事情など，複数の情報を総合して判断されている．
そのため，単純な条件での分類は困難である．
例えば，経由地の一部が同じでも，同じ方面にはできない．
終点が同じでも，どこから来たのかが違うと同じ方面にはできない．
同じ系統でも，異なる経路を持つ便が存在する．系統も完全な根拠にはならない．
往路と復路で経路が異なる場合もある．
路線網の形状が同じでも，区切り方が地域の人流や地形に左右されるため，一意に決まらない．
そこで，AIエージェントに判断基準と情報を与え，解析させる．

\subsection{本研究で解決する課題}
本研究では，以下の課題を解決する．

\begin{itemize}
    \item headsignの文字列から利用者向けの意味情報を抽出する
    \item GTFSに直接存在しない「方面」を路線網から分類する
\end{itemize}
